\section{Synthesis and focused next experiment}
\label{sec:jesper-synthesis}
\contributionauthor{Jesper Eggers}
The completed experiments support three distinct conclusions. Temporal
masking preserves more viable escapes and admits fewer doomed actions in
the generated suite. With final-network weights frozen, it improves score
and reduces self-destruction under both tested opponent populations.
Learned ranking adds further score beyond the temporal mask, but increases
self-destruction relative to the corresponding random control. These
results establish complementary contributions of masking and learned
selection under controlled deployment conditions.

The competitive trace identifies a concrete remaining weakness: an escape
can depend on a tile that an opponent reaches before movement executes.
A focused next experiment would test immediate escape destinations for
conflicts with reachable opponent occupancy, then compare the resulting
constraint with the existing mask using frozen weights and paired games.
Greater conservatism may prevent useful attacks, so score and survival
must be assessed jointly. Retraining would follow only if this deployment
test motivates it.

For future competitions, the framework should export initial-layout seeds,
requested actions, invalid-action events, model hashes, and per-agent
episode outcomes through one evaluation utility. Separate generators for
board layouts and action order would make such comparisons reproducible
by default.
